\documentclass[12pt]{article}
\usepackage[T1]{fontenc}
\usepackage[utf8]{inputenc}
\usepackage{lmodern}
\usepackage{textcomp}
\usepackage{enumitem}
\usepackage{geometry}
\geometry{margin=1in}
\begin{document}
\begin{center}
Answer Key - History - K-12 Level Assessment 13 \
\small (Answers are ordered exactly as in the question paper.)
\end{center}

\section*{Multiple Choice}
\begin{enumerate}[label=\arabic*.]
\item \textbf{Answer:} B
\\ \textit{Options:} A) Occipital; B) Acheulean; C) Oldowan; D) Stadial
\\ \textit{Note:} Homo erectus is most closely associated with the Acheulean tool-making tradition because this species is known for producing distinctive bifacial hand axes and other tools that characterize the Acheulean period, marking a significant advancement in early human technology.
\item \textbf{Answer:} C
\\ \textit{Options:} A) pit-houses; B) pueblos; C) kivas; D) courtyard groups
\\ \textit{Note:} Kivas are ceremonial structures used by Native American cultures, particularly the Pueblo peoples, for religious rituals and gatherings, making them the correct answer to the question about structures that housed religious ceremonies.
\item \textbf{Answer:} A
\\ \textit{Options:} A) the need to concentrate population away from arable lands near rivers; B) the need to construct monumental works, such as step-pyramids and temples; C) the need to develop a complex social system that would allow the construction of canals; D) the use of irrigation to produce enough food
\\ \textit{Note:} The need to concentrate population away from arable lands near rivers is incorrect as a deciding factor in the development of Mesopotamian civilization because the civilization actually thrived by utilizing the fertile land near rivers for agriculture, which supported population growth and urbanization.
\item \textbf{Answer:} C
\\ \textit{Options:} A) 2 million years; B) 4 million years; C) 7 million years; D) 15 million years
\\ \textit{Note:} The correct answer of 7 million years reflects the estimated time since humans and chimpanzees shared a common ancestor, based on genetic and fossil evidence, indicating their separate evolutionary paths.
\item \textbf{Answer:} C
\\ \textit{Options:} A) paleoanthropology; B) archaeology; C) primatology; D) linguistics
\\ \textit{Note:} Primatology is the subfield of anthropology that specifically studies primates, which are the nearest living relatives of humans, to gain insights into human evolution and behavior.
\end{enumerate}

\section*{True / False}
\begin{enumerate}[label=\arabic*.]
\setcounter{enumi}{5}
\item \textbf{Answer:} True
\\ \textit{Options:} A) True; B) False
\\ \textit{Note:} The Lapita pottery style is true because it represents a significant cultural tradition that emerged around 1600 BCE and is recognized for its distinctive designs, which were used by the ancestors of the Polynesians as they migrated and settled across the Pacific Islands, indicating a shared cultural identity and complex societal interactions.
\item \textbf{Answer:} False
\\ \textit{Options:} A) True; B) False
\\ \textit{Note:} Monk's Mound was primarily a ceremonial and political center rather than a burial site, serving as a platform for important structures and gatherings in the Cahokia civilization.
\item \textbf{Answer:} False
\\ \textit{Options:} A) True; B) False
\\ \textit{Note:} Olmec monuments are primarily understood to be ceremonial and religious structures rather than storage centers for surplus food, as they often feature elaborate carvings and align with spiritual practices.
\item \textbf{Answer:} False
\\ \textit{Options:} A) True; B) False
\\ \textit{Note:} The statement is false because the main Indus Valley cities, such as Mohenjo-Daro and Harappa, were fortified with defensive walls to protect against potential invasions and floods.
\item \textbf{Answer:} True
\\ \textit{Options:} A) True; B) False
\\ \textit{Note:} Severe drought led to water shortages and agricultural decline, which diminished the resources and labor available for construction at Maya ceremonial centers.
\end{enumerate}

\section*{Long Form Response}
\begin{enumerate}[label=\arabic*.]
\setcounter{enumi}{10}
\item \textbf{Answer:} Following the classical period around A.D. 800, the Maya civilization experienced significant changes marked by decline, relocation, resurgence, and a subsequent decline. Initially, many southern cities faced abandonment due to factors such as drought, warfare, and resource depletion. As a result, the Maya shifted their center of power northward to the Yucatán Peninsula, where they established new cities like Chichen Itza. This relocation led to a resurgence in trade, culture, and population as the Maya adapted to their new environment. However, by the end of the 10 th century, this new phase also faced decline due to similar challenges, including environmental stress and social upheaval, leading to the gradual fading of the Maya civilization.
\\ \textit{Note:} correct because it accurately outlines the key phases of Mayan civilization following A.D. 800, including the factors leading to their decline, the shift of power and population to the north, the cultural and economic revival during this relocation, and the eventual decline that followed.
\item \textbf{Answer:} Radioactive dating techniques are crucial for determining the age of rocks because they allow scientists to measure the decay of certain isotopes within those rocks. Two significant methods are argon-argon dating and potassium-argon dating. These techniques rely on the concept of half-life, which is the time it takes for half of the radioactive isotope to decay into a stable form. For instance, potassium-40 decays into argon-40 over a known half-life, enabling researchers to calculate the age of a rock sample based on the ratio of parent isotopes to daughter products. By using these methods, scientists can accurately date volcanic rocks and understand the geological history of the Earth, providing insights into events like volcanic eruptions and the formation of landforms.
\\ \textit{Note:} correct because it accurately describes how radioactive dating techniques, specifically argon-argon and potassium-argon dating, utilize the concept of half-life to measure the decay of isotopes, allowing scientists to determine the age of rocks.
\end{enumerate}

\vfill
\noindent\textit{End of Answer Key}
\end{document}